\section{Interleaved codes as codes over an extension}

\setcounter{lemma}{4}\begin{lemma}[Interleaved code $=$ extension-field code]\label{lem:ext}
Let $\E\supseteq\F$ be an extension of degree $t$, $e_0,\dots,e_{t-1}$ an $\F$-basis of
$\E$, and $\varphi:\F^t\to\E$, $a\mapsto\sum_s a_se_s$, applied pointwise to words on
$L\subseteq\F$. Then:
\begin{enumerate}
\item $\varphi$ is an $\F$-linear bijection $\RS[\F,L,d]^t\to\RS[\E,L,d]$;
\item $\varphi$ preserves agreement sets $\{x:u(x)=v(x)\}$, hence $\Delta$ and
$\agree_\mu$, for any two words;
\item $\varphi\big(\sum_j z_ju_j\big)=\sum_j z_j\,\varphi(u_j)$ for $z_0,\dots,z_l\in\F$.
\end{enumerate}
\end{lemma}

\begin{proof}
(1) $P\in\E[X]_{<d}$ decomposes uniquely as $P=\sum_s e_sP_s$ with
$P_s\in\F[X]_{<d}$ (expand each coefficient in the basis), and for $x\in L\subseteq\F$,
$P(x)=\sum_s e_sP_s(x)=\varphi\big((P_s(x))_s\big)$.
(2) $\varphi$ is injective, so $u(x)=v(x)\iff\varphi(u(x))=\varphi(v(x))$.
(3) $\varphi$ is $\F$-linear.
\end{proof}
\uses{\cref{def:codes,def:agree}}\usedby{L6, L8}

\begin{remark}
Item (3) is the only place where it matters that the challenges lie in $\F$ and not in
$\E$; it is what L6 uses.
\end{remark}
